\section{Output Interface and Decoding Diagnostics}
\paragraph{How is a location expressed?}
Output syntax and coordinate scale create a second mechanism for score changes.
Figure~\ref{fig:protocol} holds each set of raster inputs fixed across the
readout variants. Several configurations respond differently to pixel and
normalized requests. Molmo, for example, produces responses on a 0--100 scale in
some JSON/pixel prompts, while its two-line normalized prompts produce
0--1000-style points. Qwen3-VL-8B commonly retains a normalized-style output
even when asked for pixels. The raw responses reveal these scale mismatches before target-box scoring.

\begin{figure}[t]
\centering
\includegraphics[width=\linewidth]{protocol.pdf}
\caption{Readout-dependent paired score changes in the 30-pair raster archive.
Each cell uses the requested coordinate frame and all 30 queries in each
condition. The heatmap reports percentage-point differences and the accompanying coverage
summary counts parsed numeric outputs. Full counts are in
Table~\ref{tab:protocol-full}.}
\label{fig:protocol}
\end{figure}

The response format therefore changes the meaning of a measured material
effect. A parser may discard a response, a coordinate mismatch may displace a
correctly identified location, or a valid point may genuinely move on the image.
Our analysis distinguishes these cases using raw text, explicit coordinate
transforms, and pixel distances. A miss caused by a scale mismatch therefore has a different explanation from
a point that moves off the object under a fixed readout.

\paragraph{Decoding variability}
Three stochastic runs per image provide another view of spatial readout.
On the recorded raster inputs, Qwen3-VL-8B has a mean pairwise seed distance of
7.66 pixels, compared with 24.34 for Qwen3-VL-2B and 27.07 for Gemma4.
The corresponding hit-status variability is 0\%, 16.7\%, and 10.0\%.
Table~\ref{tab:sampling} gives the complete-query counts. These differences
motivate fixed decoding when comparing the two image conditions.
